\chapter{Related Work}
\label{chap:related-work}

Among the studies on AI-assisted development, several studies
are relevant to the problems addressed in this thesis. This chapter
reviews work on prompt consolidation and project handoff, errors in consolidated
requirements, requirement traceability evaluation, methods for keeping
specifications up to date, and semantic conflict detection. The final section
brings these areas together and identifies the research gap addressed by this
thesis.

\section{Prompt Consolidation and Project Handoff}
\label{sec:related-consolidating}

The work most directly related to prompt consolidation is presented by
\citet{mondal2024consolidating}. The study shows how users solve an issue through
sequences of prompts and propose combining those prompts into a single prompt.
Their main goal is to reduce repeated interaction and produce the desired
answer more efficiently. They do not study about whether every statement in the
consolidated prompt can be traced to the original prompt sequence. They also
do not examine how later prompts remove, undo, or replace earlier
requirements.

Recent work has also studied how an AI coding task can be continued after a
handoff. \citet{kc2026handoffdebt} describe \emph{handoff debt}: the work a
new coding agent must repeat to understand an interrupted task. Their focus is
on transferring repository state and task information so that another agent
can resume the work. This differs from the present thesis, which studies how
faithfully the current project intent can be represented in one master
prompt.

Another approach is to preserve the interaction history rather than
consolidate it. \citet{santoni2026cmv} propose a graph-based memory system that
keeps previous interaction turns available for later retrieval. This avoids
reducing the full history to one document. The present research studies a
different case: the history is deliberately turned into a single portable
artifact that can be read, stored, and handed over.

\section{Errors in Consolidated Requirements}
\label{sec:related-loss}

The main failure studied in this thesis is that a generated master prompt may
contain requirements that are not supported by the user's prompt history.
This can be understood through existing work on hallucination.

\citet{ji2023survey} distinguish between intrinsic and extrinsic
hallucination. Intrinsic hallucination occurs when generated content
contradicts its source. Extrinsic hallucination occurs when generated content
cannot be supported by the source. This, by definition, is relevant to master-prompt
consolidation. A technology, filename, validation rule, or design choice that
the user never requested is unsupported requirements and is therefore close to
extrinsic hallucination. 

A different error occurs when the master prompt retains an old requirement
after a later prompt has changed, removed, or replaced it. This content
originally appeared in the prompt history, so it is not unsupported in the
same way as extrinsic hallucination. Instead, it reflects a failure to 
follow the order of user's prompts and recover the user's current intent.

Related work has also used language models to obtain software requirements
from conversations. \citet{voria2025recover} propose RECOVER, which finds
and generates system requirements from stakeholder conversations. This shows
that unstructured dialogue can be used as a source of requirements. However,
RECOVER focuses on generating requirements from stakeholder discussions. Unlike 
the problem addressed by \citet{voria2025recover}, the prompt history already
contains direct development instructions, but their current state must be 
recovered after modifications, removals, replacements, and changes of direction.

\section{Evaluating Requirement Traceability}
\label{sec:related-metrics}

Word similarity alone is not reliable enough for evaluating a master prompt. Two
requirements may use almost the same words while conflicting with each other in meaning.
In contrast, two requirements can use different words while expressing the same
requirement. For this reason, this study uses entailment-based evaluation to check whether the meaning of a candidate requirement is supported by the ground truth.

\citet{zha2023alignscore} propose AlignScore, a factual-consistency metric that
measures whether a claim is supported by a context. The model treats the two
inputs differently: one is the source context and the other is the claim being
checked. In this thesis, the manually constructed ground truth is used as the
context and each candidate requirement is used as the claim:

\[
\mathrm{ALIGN}
(\mathrm{context}=\mathrm{ground\ truth},\
\mathrm{claim}=\mathrm{candidate\ requirement}).
\]

A low score means that the candidate requirement is not well supported by the
ground truth. This makes AlignScore suitable for finding agent-added details
and outdated requirements.

The evaluation is performed at the requirement level rather than on the full
master prompt. This choice is related to FActScore and Claimify.
\citet{min2023factscore} evaluate long-form text by first separating it into
atomic facts and then checking whether each fact is supported.
\citet{metropolitansky2025claimify} study how complex text can be separated
into clear and independent claims. This thesis uses the same
general idea but treats each unit as an atomic software requirement.

This requirement-level evaluation gives each requirement its own score, making it
possible to identify exactly which parts of a master prompt are unsupported.
It also prevents well-supported requirements from hiding a small number of
incorrect ones. The scores from individual requirements
can then be used to calculate the untraceable rate and to support 
human review of the errors found in each candidate.

Traceability has also been studied between other software artifacts.
\citet{alor2025traceability} evaluate whether language models can create trace
links between software documentation and source code. Their artifacts and
evaluation setting are different from those in this thesis, but there is 
common concern about whether one artifact is supported by another.
However, in this thesis, each generated requirement is checked against
the ground-truth master prompt to determine whether 
it is supported by the user's original requirements.

\section{Keeping an Up-to-Date Specification}
\label{sec:related-spec}

Spec-driven development provides a related view of how natural language
requirements can be used during AI-assisted development. \citet{piskala2026specdriven} describe approaches in which a specification is
kept as the main description of what the software should do, while code is
generated or checked against it.

The \texttt{spec\_guided} method in this thesis uses a similar principle.
Instead of asking an agent to rebuild the current requirements from the full
prompt history only at the end, it keeps a specification updated throughout
the session. The specification records the requirements that are currently
valid, while the original prompts are kept separately as their source.

This thesis does not treat the specification as a complete replacement for
code, nor does it claim to implement a full spec-driven development system.
It uses the maintained-specification idea for a narrower purpose: reducing
non-traceable content during master-prompt construction and giving the agent
a current document against which later instructions can be checked.

Prompts can also become maintained artifacts in software projects.
\citet{tafreshipour2025wild} study how prompts stored in software repositories
change over time and report inconsistencies between prompt revisions. Their
work examines repository-based prompts rather than the developer--agent prompt
histories studied in this thesis. However, it supports the view that prompts
may need to be stored and maintained carefully when they remain part of a
project over time.

\section{Semantic Conflict Detection}
\label{sec:related-vc}

A text document can merge successfully while still containing requirements
that conflict in meaning. Such a conflict is called a semantic conflict.
Related work in collaborative code generation shows the same difference 
between text-level consistency and semantic correctness.

\citet{pugachev2025codecrdt} propose CodeCRDT, which allows multiple coding
agents to edit shared code while avoiding text-level merge conflicts.
Their preliminary inspection suggests that 5--10\% of the inspected runs
still contain semantic problems such as duplicate declarations and type
mismatches. This result is not a complete measurement of semantic-conflict
frequency, but it shows an important point: a clean text-level merge
does not guarantee a correct result.

CodeCRDT studies conflicts in source code. Unlike source code, it cannot be checked by a compiler, type checker, or test suite, and there is no exact score that can always determine whether two requirements conflict. The conflict must therefore be identified by comparing their meanings. Thus, this thesis shows that the conflict must be found 
by comparing the meaning of the incoming instruction with the meaning of 
the existing specification.

The work most closely related to the conflict part of this thesis is
\emph{Semantic Commit} by \citet{vaithilingam2025semanticcommit}. Their system
helps users add new information to an existing intent specification, such as a
Cursor rules file or game design document. It uses a knowledge graph and
retrieval to find existing information that may be affected by the update.
The retrieved information is then classified as a direct conflict, an
ambiguous conflict, or a non-conflict. The interface presents the result to
the user and supports human review and resolution.

Semantic Commit also evaluates a baseline called DropAllDocs, which provides
the complete existing information store to the model instead of retrieving
selected parts. The present thesis uses a related full-context idea because
the tested master prompts are small enough to fit in the model context.
However, it does not reproduce Semantic Commit's knowledge-graph retrieval
pipeline or its user interface.

The main difference is the starting point. Semantic Commit assumes that an
intent specification already exists and studies how new information should be
added to it. This thesis first studies whether such a specification can be
constructed faithfully from a long prompt history. It then reuses the
maintained specification to check next instructions for conflict. In other
words, Semantic Commit mainly addresses updates to an existing intent
specification. It does not study how that specification was created from
the earlier development history, while RQ1 and RQ2 examine how a traceable master prompt is
created before updating the specification.

The decision not to resolve a conflict automatically is also related to the
intent gap described by \citet{lahiri2026intent}. A model may produce an
implementation that appears reasonable while still differing from what the
user intended. When two requirements cannot both be true and the user's
preferred value is not clear, silently selecting one value may create the
same problem. The safer action is to report the conflict and ask the user to
decide.

\section{Research Gap}
\label{sec:positioning}

The reviewed work can be compared along three parts of the problem: producing
one consolidated artifact, checking whether its content is supported by the
source prompts, and reporting semantic conflicts when the artifact is updated.

Prompt consolidation by \citet{mondal2024consolidating} addresses the first
part, but its purpose is interaction efficiency rather than requirement
traceability. AlignScore \citep{zha2023alignscore}, FActScore
\citep{min2023factscore}, and Claimify \citep{metropolitansky2025claimify}
provide useful ideas for evaluating support at a fine level, but they do not
construct or maintain a master prompt. Spec-driven development
\citep{piskala2026specdriven} supports the use of a maintained specification,
but does not directly evaluate whether an AI-generated master prompt contains
requirements that the user never requested. CodeCRDT
\citep{pugachev2025codecrdt} shows that text-level convergence does not
prevent semantic errors, but it works on source code rather than
natural-language requirements.

Semantic Commit \citep{vaithilingam2025semanticcommit} is the closest work for
conflict detection because it checks new information against an existing
intent artifact and keeps the user involved in resolution. However, it begins
with an intent specification that has already been prepared. It does not
study whether consolidating a raw, ordered prompt history introduces
unsupported requirements or retains values that were later changed.

Among the works reviewed in this chapter, the specific combination studied in
this thesis remains open: evaluating the traceability of a master prompt built
from a chronological AI coding history, comparing single-pass consolidation
with a maintained-specification method, and testing whether the same method
can report later semantic conflicts without silently resolving them.
